
\noindent 
Let $\sigma$ be a skew Young diagram (Definition~\ref{def:young}).
The main result of this paper is a new formula for the \emph{skew stable Grothendieck polynomial} $G_\sigma$ of Lascoux--Sch{\"u}tzenberger and Fomin--Kirillov~\cite{fomin-kirillov-grothendieck,lascoux-schutzenberger} as a linear combination of skew Schur functions $s_\lambda$ on related shapes $\lambda$. As was demonstrated by Buch, the coefficients of $G_\sigma$ have a combinatorial interpretation in terms of set-valued tableaux~\cite{buch-littlewood-richardson}, and
our original motivation for this paper came from a recent geometric result, proved in a companion paper~\cite{chan-pflueger-euler}, identifying Euler characteristics of Brill--Noether varieties up to sign as counts of set-valued standard tableaux. 

It is natural to ask for a linear expansion of $G_\sigma$ in terms of other symmetric functions, particularly the basis of Schur functions. 
Such an expansion was obtained by Fomin--Greene, who in fact obtained such expansions for a wide class of symmetric functions including stable Grothendieck polynomials associated to arbitrary permutations~\cite{fomin-greene-noncommutative}. (Note that the stable Grothendieck polynomials of $321$-avoiding permutations precisely correspond to stable Grothendieck polynomials of skew shapes as in~\cite{buch-littlewood-richardson}, by a theorem of Billey--Jockusch--Stanley~\cite{billey-jockusch-stanley-some}.) Buch's expansion of skew Grothendieck polynomials in terms of Grothendieck polynomials of \emph{straight} shapes, along with Lenart's expansion of the latter into Schur functions, provides another route to such an expansion~\cite{buch-littlewood-richardson,lenart-combinatorial}. 

Our main theorem expresses $G_\sigma$ instead as a linear combination of \emph{skew} Schur functions $s_\lambda$. 
The coefficients of the linear combination have explicit combinatorial interpretations which we provide; they count appropriate auxiliary tableaux. We state the result below, postponing all definitions to the next section.

\begin{theo}\label{thm:comb}
For any connected skew shape $\sigma$, the skew Grothendieck polynomial $G_\sigma$ admits a linear expansion 
\[
G_\sigma = \sum_{\mu\supseteq \sigma} (-1)^{|B(\mu/\sigma)|} a_{\sigma,\mu} \cdot s_\mu
\]
where the $a_{\sigma, \mu}$ are nonnegative integers and the $s_\mu$ are skew Schur functions. Here $a_{\sigma,\mu}$ is the product of the following two integers:


\begin{enumerate}
\item\label{theo1.1_1} the number of row-weakly-bounded semistandard tableaux of shape $A(\mu/\sigma)$, and
\item\label{theo1.1_2} the number of row-bounded, reverse row-strict tableaux of shape $B(\mu/\sigma)$.
\end{enumerate}
\end{theo}

\noindent Here $A(\mu/\sigma)$ and $B(\mu/\sigma)$ are the subshapes of $\mu$ lying above and below $\sigma$, respectively. In fact, this statement is a specialization of a more general formula for row-refined skew stable Grothendieck polynomials that we obtain in Theorem~\ref{t:refinedcount}. 

Theorem~\ref{thm:comb} generalizes Lenart's theorem from 2000 expanding Grothendieck polynomials for non-skew shapes into non-skew Schur functions~\cite{lenart-combinatorial}; in fact, that result is a visible specialization. We explain this connection in detail in Remark~\ref{rem:connection-to-lenart}. We also give a theorem in the other direction, Theorem~\ref{t:back}, expressing $s_\sigma$ in terms of polynomials $G_\mu$, for skew shapes $\sigma$. This generalizes an analogous theorem of Lenart from the same paper.



To be clear, skew Schur functions, since they include Schur functions properly, are evidently not a basis for the space of symmetric functions; thus the coefficients of our expansion are not canonical. 
To provide a point of comparison, a result in the literature that is similar in spirit to Theorem~\ref{thm:comb} is the skew Pieri rule of Assaf--McNamara, in which the product of a skew shape and a rectangle is expressed in terms of other skew shapes~\cite[Theorem~3.2]{assaf-mcnamara-pieri}. Again, this expression is necessarily noncanonical, but it is combinatorially natural using an insertion algorithm. Our proof also uses a new insertion algorithm for skew set-valued semistandard tableaux that is related to previous work of Bandlow--Morse, and indeed our algorithm may be interpreted as extending to the skew case some of their results~\cite[\S~5]{bandlow-morse-combinatorial}. In fact, it recalls earlier work of Sagan--Stanley row insertion for skew (non-set-valued) tableaux~\cite{sagan-stanley-robinson}, as well as Buch's ``uncrowding'' algorithm on set-valued tableaux~\cite[\S~6]{buch-littlewood-richardson}. We also note that using insertion operations to derive such combinatorial identities has been carried out previously, in the form of Hecke insertion operations studied in~\cite{bksty-stable}.

%\medskip

\subsection*{Motivation from geometry} 
Our original motivation came from a recent result in Brill--Noether theory that we prove in a pair of companion papers (\cite{chan-pflueger-euler}, together with~\cite{chan-pflueger-relative} which proves an auxiliary result). 
\begin{theo}[\cite{chan-pflueger-euler}] Fix $r,d\ge 0$ and nondecreasing sequences $\alpha,\beta\in\ZZ^{r+1}_{\ge 0}$. Let $(X,p,q)$ be a general twice-pointed curve of genus $g$ over an algebraically closed field.
Then the algebraic Euler characteristic of the Brill--Noether variety $\Grdab(X,p,q)$ is 
\begin{multline*}
\chi(\Grdab(X,p,q))\\ = (-1)^{g-|\sigma|} \cdot \#(\text{standard set-valued tableaux on $\sigma$ of %} %\\ &\text{
content }\{1,\ldots,g\}).
\end{multline*}
\end{theo}	 
\noindent Here $\sigma$ is the skew-shape obtained from an $(r+1)\times (g-d+r)$ rectangle by adding $\alpha_r \le \cdots \le \alpha_0$ boxes down the left side and $\beta_0\ge \cdots \ge \beta_r$ boxes down the right side. The partitions $\alpha$ and $\beta$ encode some ramification conditions imposed at the two marked points $p,q$ of $X$. Roughly speaking, from the geometric perspective it is natural to seek formulae for set-valued tableaux in terms of skew Schur functions, which do not give preference to one marked point over the other, rather than (straight) Schur functions, which do.
Indeed, 
our result provides a combinatorial explanation of the main theorem of~\cite{anderson-chen-tarasca-k-classes}, as we shall explain further in Remark~\ref{rem:connection-to-act}. It also generalizes a theorem with L\'opez and Teixidor i Bigas which computes genera of Brill--Noether curves~\cite{clpt}. (That case corresponds to the situation in which there is exactly one more label than the number of boxes.) 
In addition, a recent result of Reiner--Tenner--Yong 
is also a special case of Theorem~\ref{thm:comb}, and in fact, their work inspired some of the results here~\cite[Corollary~3.11]{reiner-tenner-yong}. 

\section{Preliminaries}\label{sec:prelim}

We now give some preliminaries on tableaux. First we define a skew Young diagram: this is almost the same as the usual definition, except that we only record the set of boxes in a diagram rather than remembering a formal difference of two partitions.
Fix the partial order $\preceq$ on $\ZZ^2$ given by $(x,y)\preceq (x',y')$ if $x\le x'$ and $y\le y'$. 

\begin{defi}\label{def:young}\ \\*[-1.2em]
\begin{enumerate}
\item\label{defi2.1_1} 
A skew Young diagram is a finite subset $\sigma \subset \ZZ_{>0}^2$ that is closed under taking intervals. In other words, $\sigma$ has the property that if $(x,y)$ and $(x',y')\in \sigma$ with $(x,y)\preceq (x',y')$, then 
\begin{center}
$\{(x'',y''):(x,y)\preceq (x'',y'') \preceq (x',y')\}\subseteq \sigma.$
\end{center}
\item\label{defi2.1_2} A skew Young diagram is called a \emph{Young diagram} if $\sigma$ is empty or has a unique minimal element.
\end{enumerate}
\end{defi}

Skew Young diagrams are sometimes also called \emph{skew shapes}, and skew Young diagrams having a unique minimal element will sometimes be called \emph{straight shapes} for emphasis. 
In accordance with the English notation for Young diagrams, we will draw the points of $\ZZ^2$ arranged with $x$-coordinate increasing from left to right, and $y$-coordinate \emph{increasing} from top to bottom, \eg
\[
\begin{array}{ccc}
(1,1) & (1,2) & \cdots \\
(2,1) & (2,2) & \\
\vdots & &
\end{array}
\]
Furthermore, we will draw, and refer to, the members of $\sigma$ as boxes, as usual, and we let $|\sigma|$ denote the number of boxes in $\sigma$. We shall assume throughout for convenience that $\sigma$ is a connected shape, \ie its Hasse diagram is connected (see Remark~\ref{rem:disconn} on the disconnected case). 

\begin{defi} \label{def:tableau}
A \emph{tableau} of shape $\sigma$ is an assignment $T$ of a positive integer, called a \emph{label}, to each box of $\sigma$. \begin{enumerate}
\item\label{defi2.2_1} A tableau $T$ of shape $\sigma$ is \emph{semistandard} if the rows of $\sigma$ are weakly increasing from left to right, and the columns of $\sigma$ are strictly increasing from top to bottom.
\item\label{defi2.2_2} A tableau $T$ of shape $\sigma$ is \emph{standard} if it is semistandard and furthermore each integer $1,\ldots,|\sigma|$ occurs exactly once as a label.
\end{enumerate}

\end{defi}

\begin{defi}[\cite{buch-littlewood-richardson}] \label{def:svt}
A \emph{set-valued tableau} of shape $\sigma$ is an assignment of a nonempty finite set of positive integers to each box of $\sigma$. 

Given sets $S,T\subseteq \ZZ_{>0}$, we write $S<T$ if $\max(S)<\min(T)$, and we write $S\le T$ if $\max(S)\le\min(T)$. Then we extend the definitions of \emph{semistandard} and \emph{standard} tableaux to set-valued tableaux.
\begin{enumerate}
\item \label{it:sstd}A set-valued tableau $T$ of shape $\sigma$ is \emph{semistandard} if the rows of $\sigma$ are weakly increasing from left to right, and the columns of $\sigma$ are strictly increasing from top to bottom.
\item \label{it:std} A set-valued tableau $T$ of shape $\sigma$ is \emph{standard} if it is semistandard and furthermore the labels are pairwise disjoint sets with union $\{1,\ldots,r\}$ for some $r \ge |\sigma|$.
\end{enumerate}
Denote by $\Semi(\sigma)$ the set of all semistandard set-valued tableaux on $\sigma$.
\end{defi}


Let $\bfc=(c_1,c_2,\ldots)$ be a nonnegative integer sequence that is eventually zero. We say that a tableau or set-valued tableau $T$ of shape $\sigma$ has content $\bfc = c(T)$ if label $i$ appears exactly $c_i$ times, for all $i$. 
Write $|T| = |c(T)| = \sum c_i$ for the total number of labels.

\begin{defi} \label{def:schur-groth}\ \\*[-1.2em]
\begin{enumerate}
\item\label{defi2.4_1}
For any skew shape $\sigma$, the \emph{skew Schur function} $s_\sigma$ is
\[
s_\sigma =\sum_T {\mathbf{x}}^{c(T)}
\]
as $T$ ranges over all semistandard fillings of $\sigma.$
\item\label{defi2.4_2}
For any skew shape $\sigma$, the \emph{skew stable Grothendieck polynomial} $G_\sigma$ is
\[
G_\sigma =\sum_T (-1)^{|T|-|\sigma|}{\mathbf{x}}^{c(T)}
\]
as $T$ ranges over all semistandard set-valued fillings of $\sigma.$
\end{enumerate}
\end{defi}

Given a set-valued tableau $T$ of shape $\sigma$, define the \emph{excess} of $T$, denoted $e(T)$, as the vector $\bfe = (e_1,e_2,\ldots)$ in which $e_i$ records the number of labels in row $i$ in excess of the number of boxes in row $i$. Therefore $|\sigma|+|e(T)|=|c(T)|$. 

\begin{rema} \label{rem:disconn}
Definition~\ref{def:schur-groth} does not require $\sigma$ to be a connected skew shape, but there is little loss of generality in focusing on the connected case, as we do in this paper. If $\sigma$ is a union of two disconnected parts $\sigma_1, \sigma_2$, then a filling $T$ of $\sigma$ is semistandard if and only if the resulting fillings $T_1$ of $\sigma_1$ and $T_2$ of $\sigma_2$ are semistandard, since no box in $\sigma_1$ is comparable by $\preceq$ to a box of $\sigma_2$. Also, $c(T) = c(T_1) + c(T_2)$. It follows that $G_\sigma = G_{\sigma_1} G_{\sigma_2}$. Hence Grothendieck polynomials of disconnected skew shapes factor into those of connected skew shapes.
\end{rema}

\section{Row insertion for skew set-valued tableaux}

We now set notation for a refinement of the Grothendieck polynomial based on the excess statistic, and we prove a theorem expressing it linearly in terms of skew Schur functions. Note that the idea of introducing an additional parameter into the Grothendieck polynomial goes back already to~\cite{fomin-kirillov-grothendieck}.

\begin{defi} \label{def:rg} Let $\sigma$ be a skew Young diagram.
We define the \emph{row-refined skew stable Grothendieck polynomial} of $\sigma$ to be the power series 
\[
RG_\sigma(\mathbf{x};\mathbf{w}) = \sum_{T\in \Semi(\sigma)} (-1)^{|e(T)|} \mathbf{x}^{c(T)} \mathbf{w}^{e(T)}.
\]
\end{defi}
\noindent Thus $RG_\sigma(\mathbf{x};\mathbf{1}) = G_\sigma(\mathbf{x})$, so the usual skew stable Grothendieck polynomial is obtained as a specialization.


\begin{defi}Let $\mu$ be a skew Young diagram.
\begin{enumerate}
\item\label{defi3.2_1} A tableau $T$ of shape $\mu$ is \emph{reverse row-strict} if its rows are strictly decreasing from left to right, and its columns are weakly decreasing from top to bottom. 
\item\label{defi3.2_2} A tableau $T$ of shape $\mu$ is \emph{row-bounded}, respectively \emph{row weakly-bounded}, if for every box $(i,j)$ in $\mu$, $T(i,j) < i$, respectively $T(i,j)\le i$.
\end{enumerate}
\end{defi}

We henceforth adopt the following convention governing containment of Young diagrams. 
\begin{enonce}{Convention}\label{conv}
Fix $\sigma$ a skew shape; we take the numbering of the rows of $\sigma$ to start at $1$ at the top. For another skew shape $\lambda$, we write $\lambda \supseteq \sigma$ if every box of $\sigma$ is a box of $\lambda$, \emph{and furthermore every box of $\lambda$ is in the same column as some box of $\sigma$}. In other words, we will only consider skew shapes $\lambda\supseteq \sigma$ that occupy the same set of columns as $\sigma$. They are not allowed to extend $\sigma$ to the right or to the left.
\end{enonce}

By Convention~\ref{conv}, if $\sigma$ is a connected skew shape and $\lambda\supseteq \sigma$, then $\lambda - \sigma$ consists of a set of boxes above $\sigma$ and a set of boxes below $\sigma$. Write $A(\lambda/\sigma)$ and $B(\lambda/\sigma)$ for these respective skew Young diagrams; $A$ and $B$ stand for \emph{above} and \emph{below.} We emphasize that, contrary to some conventions, $\lambda$ may extend $\sigma$ both above and below.



